\section{A further reduction to 226 states}
\label{sec:extension226}

The calibrated 227-state layout admits one additional closed source
rotation and a preceding-bit fold.  The final table has 226 working
states and 452 transition rows.  This extension changes or deletes nine
rows of the inherited 227-state table.

\subsection{The source rotation}

At source prefix \texttt{11101001001110011100}, replace
$((A,B),C)$ by $(A,(B,C))$.  The three chunks contain 28, 81, and 138
instruction occurrences, respectively, and all have boundary signature
$(0,0,\mathrm{false})$.  The corresponding decoder states are the outer
state 210 and its left child 177.  State 177 has no other incoming
transition before the rotation.  The four edited rows are
\[
\begin{array}{c|cc}
 &\text{before}&\text{after}\\\hline
177/0&0\ R\ 165&0\ R\ 240\\
177/1&1\ R\ 240&1\ R\ 215\\
210/0&0\ R\ 177&0\ R\ 165\\
210/1&1\ R\ 215&1\ R\ 177.
\end{array}
\]
The same closed-chunk argument applies as in
Section~\ref{sec:source227}.  Both independent source resolvers check
all 5,121 ordinal operations and 6,656 guarded cases, together with the
no-stray successor conditions.  The source still has 211 distinct DAG
nodes, maximum depth 37, and maximum zero-run six.

\subsection{The four roles in the fold}

The discovery tuple $(86,508,182,177)$ means: retain state 86, whose
entries have a preceding zero; delete state 508, whose entries have a
preceding one; use state 182 as the backward-read helper when the
current bit is zero; and use state 177 as that helper when the current
bit is one.  It is not a four-state deletion.

Before the fold the two decoder states are
\[
\begin{array}{c|cc}
 &0&1\\\hline
86&0\ R\ 98&1\ R\ 99\\
508&0\ R\ 236&1\ R\ 233.
\end{array}
\]
State 86 is entered only from $61/0$, $154/0$, and $233/0$, all writing
zero and moving right.  State 508 is entered only from $242/1$, which
writes one and moves right.  Neither is initial.  The local preceding
bit therefore records which of these two former states was entered.

Delete state 508, redirect $242/1$ to state 86, and set
\[
 86/0=0\ L\ 182,\qquad 86/1=1\ L\ 177.
\]
The following existing rows complete the routing:
\[
\begin{array}{c|cc}
\text{helper}&\text{preceding zero}&\text{preceding one}\\\hline
182&0\ R\ 187&1\ R\ 212\\
177&0\ R\ 240&1\ R\ 215
\end{array}
\qquad
\begin{array}{c|c}
187/0&0\ R\ 98\\
212/0&0\ R\ 236\\
240/1&1\ R\ 99\\
215/1&1\ R\ 233.
\end{array}
\]
For each of the four combinations of preceding and current bit, one
old transition is replaced by three new transitions.  Both executions
end in the same state, with head displacement $+1$, and preserve the
whole tape.  Only the preceding and current cells are read.  The fold
checker executes all four blocks with arbitrary tape outside those
cells and checks every ordinary transition and incoming port.  This
gives a positive-block simulation from the common start state on every
binary tape, between the rotated 227-state table and the 226-state
table.  It preserves finite halting and infinite runs.

\subsection{Reconstruction and startup}

The portable \texttt{reconstruct226.py} script pins the 227-state input,
replays the single source rotation and the exact fold, checks the two
source graphs, and writes the final source, table, and block receipts.
The final inherited table has SHA-256
\begin{quote}\small\ttfamily
30082fe31158b8ae049a1bf15f56acff\\
aae71f9031642e6fa7ab75555e466be4
\end{quote}
where the displayed line break is not part of the digest.

No changed or deleted row is visited during the complete blank startup.
Both tables therefore reach exactly the same dispatcher configuration
after 409,835 transitions: state 96, head position $-5$, the same 351
marked cells, and zeros everywhere else.  The final unbounded binary
runtime certificate is applied directly to the 226-state table and its
source, in addition to these source, local-fold, and startup checks.
